\section{Otimização de Portfólio de Contratos de Energia}
\label{sec:rl_portfolio}

A literatura sobre gestão de portfólio no setor elétrico brasileiro é predominantemente inspirada na teoria de Markowitz, transposta do mercado financeiro para o contexto de contratos de energia. Nessa abordagem, o problema consiste em determinar a alocação ótima entre diferentes tipos de contratos, incluindo bilaterais de prazo fixo, contratos no mercado \emph{spot} e, em alguns casos, contratos com fontes renováveis, de modo a maximizar o retorno esperado para um dado nível de risco.

\citeonline{gunn2012} aplica o modelo de Markowitz à comercialização de energia proveniente de novos empreendimentos, como biomassa, eólica e pequenas centrais hidrelétricas, no mercado brasileiro, modelando as incertezas de oferta e as regras de mercado por meio de simulações que alimentam a fronteira eficiente risco-retorno. \citeonline{guerreiro2017} avança nessa direção ao questionar a adequação da variância como medida de risco para o setor elétrico: dado que a distribuição do PLD é assimétrica e de cauda pesada, penalizar simetricamente ganhos e perdas extremas distorce a decisão de contratação. O autor propõe o CVaR como medida mais coerente, aplicando-o à otimização de portfólios de comercialização no ACL brasileiro. \cite{munhoz2018} segue linha semelhante, analisando o nível ótimo de contratação a preço fixo versus exposição ao PLD para um consumidor livre, com minimização de custo e risco via Markowitz.

Uma limitação comum a esses três trabalhos é a formulação estática do problema: o portfólio é otimizado em um único estágio, com base em uma distribuição de cenários gerada \emph{a priori}, sem modelagem explícita da sequência de decisões ao longo do tempo. Essa simplificação ignora o fato de que contratos bilaterais têm duração superior a um mês, comprometendo recursos financeiros e posições físicas em estágios futuros, e que novas informações sobre o PLD são reveladas progressivamente, permitindo ao agente adaptar sua estratégia de contratação.

\citeonline{josemaria2018} e \cite{carvalho2021} abordam o problema sob uma perspectiva mais descritiva, detalhando os tipos de contratos disponíveis no ACL, os riscos associados a cada modalidade e as estratégias adotadas por agentes na prática. \cite{carvalho2021} destaca que a literatura acadêmica tende a focar em modelos de previsão de preços, mas falha ao não incorporar explicitamente o perfil de aversão ao risco dos agentes reais, observação que reforça a necessidade de modelos que tratem a incerteza de forma estruturada.

O trabalho de \citeonline{pedrini2019} representa o avanço mais próximo da abordagem adotada neste trabalho dentro da literatura nacional: a autora formula a decisão de contratação de um consumidor livre como um problema estocástico de dois estágios, com a possibilidade de contratar geração eólica, e resolve o modelo por decomposição de Benders. A formulação de dois estágios captura a distinção entre decisões tomadas antes e após a revelação das incertezas, incorporando a não antecipatividade de forma estrutural. No entanto, o horizonte de dois estágios é uma simplificação para o problema de uma comercializadora, cujos contratos têm duração de vários meses e cujas decisões de contratação se repetem mensalmente ao longo de todo o ano.

Em síntese, a literatura nacional sobre portfólio de contratos de energia avançou de formulações estáticas baseadas em média-variância para abordagens com medidas de risco mais adequadas à assimetria do PLD e, num caso, para a estrutura de dois estágios. No entanto, formulações que combinem decisão sequencial ao longo de múltiplos meses com não antecipatividade explícita e restrições financeiras ao longo de todo o horizonte, no contexto da comercializadora no ACL brasileiro, ainda carecem de representação na literatura.
